% !TeX root = ../../Book2.tex
\section{多线性积分形式}
\label{b2:sec:2501}

多线性估计保留多个输入之间的相互作用。当几个频率片的法向彼此分离时，乘积可能比单个延拓函数有更好的可积性；若相位可分别吸收到各输入中，则再大的参数也没有统一振荡增益。

\subsection{多线性型}
\label{b2:subsec:250101}

一个基本模型为
\[
 \Lambda_\lambda(f_1,\ldots,f_k)
 =\int_{\mathbb R^d}e^{2\pi\mathrm i\lambda P(x)}a(x)
                 \prod_{j=1}^k f_j(\pi_jx)\,\mathrm dx,
\]
其中 \(\pi_j\) 是固定线性投影，\(a\) 是紧支撑光滑振幅。讨论前应先用 Hölder 或具体投影几何确认积分有意义，再问是否有 \(\lambda^{-\epsilon}\) 的增益。
输出是多线性标量型，与\chapref{b2:ch:21}一个输入一个输出的算子不同。改变投影维数，会改变可允许的输入指数。

\ProseParagraphBreak
若投影相同，所有输入可先相乘，问题退化为单输入积分；若投影互补，输入的不同方向会限制相位抵消。多线性估计因而不仅依赖相位 Hesse 矩阵，还依赖它与整个投影族的关系。

\subsection{相位相互作用}
\label{b2:subsec:250102}

若相位可以写成 \(P=\sum_jP_j\circ\pi_j\)，则取
\(f_j(y)=e^{-2\pi\mathrm i\lambda P_j(y)}g_j(y)\) 就完全消去振荡，且各输入范数不变。若某些固定 \(g_j\) 使非振荡积分非零，就不可能有对所有输入一致的正幂次参数衰减。这称为相对于投影族的相位退化。

\ProseParagraphBreak
例如 \(\pi_1(x,y)=x\)、\(\pi_2(x,y)=y\)，相位 \(x^2+y^2\) 在这个意义下退化，虽然作为二维标量相位的 Hesse 矩阵 可逆。相位 \(xy\) 则不能写成纯 \(x\) 项加纯 \(y\) 项；其混合导数非零，并与\chapref{b2:ch:21}的 Fourier 型算子相联系。这个例子直接区分标量非退化与多线性非退化。

\subsection{缩放}
\label{b2:subsec:250103}

对无振幅的全空间齐次模型，令 \(f_j(y)\) 缩放为 \(f_j(ry)\)。左侧积分产生 \(r^{-d}\)，输入范数的乘积产生
\(r^{-\sum_j d_j/p_j}\)。尺度不变界因此要求
\(\sum_jd_j/p_j=d\)，其中 \(d_j\) 是第 \(j\) 个输入空间维数。
有固定紧支撑振幅时，尺度条件未必作为必要等式出现，但局部放大仍能检验极端集中。

\ProseParagraphBreak
抛物面还具有非各向同性缩放。切向频率缩小 \(\delta\)，法向频率变化 \(\delta^2\)，空间变量分别放大 \(\delta^{-1}\) 与 \(\delta^{-2}\)。多线性管束与解耦中的尺度都来自这一关系；在频率帽与空间球之间来回换算时，应始终保留平方关系。

\subsection{近似正交性}
\label{b2:subsec:250104}

若函数 \(u_j\) 的 Fourier 支撑有限重叠，则 Plancherel 给出
\(\|\sum_j u_j\|_2^2\le L\sum_j\|u_j\|_2^2\)。
对乘积 \(u_jv_k\)，Fourier 支撑位于两支撑的 Minkowski 和；即使原支撑互不相交，这些和仍可能重叠。因此多线性近似正交应核对输出支撑的重叠数，不能只核对各输入支撑。

\ProseParagraphBreak
在局部空间球上，乘光滑截断又会使 Fourier 支撑稍微膨胀。于是通常使用快速衰减权与有限扩大后的重叠数，而不是严格的全空间支撑等式。随机符号与平方函数提供另一种组合工具，但也要保持相关的算子常数一致。
